\documentclass[french]{book}

\usepackage[utf8]{inputenc}
\usepackage[T1]{fontenc}
\usepackage{lmodern}
\usepackage{geometry}
\usepackage{graphicx}
\usepackage{babel}
\usepackage{amsmath}
\usepackage{amsthm}
\usepackage{amssymb}
\usepackage{hyperref}
\usepackage{stmaryrd}
\usepackage{enumitem}
\usepackage{tcolorbox}
\usepackage{dsfont}
\usepackage{multirow}
\usepackage{etoc}
\usepackage{titlesec}
\usepackage{esint}
\usepackage{cancel}
\usepackage{mathdots}

\setlist[itemize]{label=$\bullet$}


\renewcommand{\P}{\mathbb{P}}
\newcommand{\N}{\mathbb{N}}
\newcommand{\Z}{\mathbb{Z}}
\newcommand{\Q}{\mathbb{Q}}
\newcommand{\R}{\mathbb{R}}
\newcommand{\C}{\mathbb{C}}
\newcommand{\K}{\mathbb{K}}
\renewcommand{\L}{\mathbb{L}}
\newcommand{\U}{\mathbb{U}}
\newcommand{\st}{^*}
\newcommand{\tq}{\middle|}
\newcommand{\inv}{^{-1}}
\newcommand{\E}{\mathbb{E}}
\newcommand{\F}{\mathbb{F}}
\newcommand{\V}{\mathbb{V}}
\renewcommand{\ker}{\text{Ker}}
\newcommand{\im}{\text{Im}}
\newcommand{\card}{\text{Card}}
\newcommand{\Tr}{\text{Tr}}

\theoremstyle{definition}
\newtheorem{ex}{Exercice}[chapter]
\newtheorem{pro}{Problème}[chapter]

\title{Exercices intéressants}

\begin{document}

\chapter{Polynômes et fractions rationnelles}

\section{Problèmes}

\begin{pro}
	Pour une fraction rationnelle non nulle $F\in\C(X)$, on note $n(F)$ le nombre de ses racines (distinctes), et $p(F)$ le nombre de ses pôles.
	\begin{enumerate}
		\item Soit $F\in\C(X)$ non constante. Montrer que $n(F)+n(F-1)+p(F)>a$, où $a$ désigne le degré du numérateur de $F$ dans son écriture sous forme irréductible. On pourra commencer par le cas où $F$ est un polynôme.
		\item Soit $P$ et $Q$ deux polynômes non constants dans $\C[X]$ tels que $P\inv(\{0\})=Q\inv(\{0\})$ et $P\inv(\{1\})=Q\inv(\{1\})$. Montrer que $P=Q$.
		\item Théorème de Mason : soit $A$, $B$ et $C$ trois polynômes non nuls premiers entre eux dans leur ensemble, non tous constants et tels que $A+B+C=0$. On note $n$ le maximum des degrés de $A$, $B$ et $C$. Montrer que $ABC$ possède au moins $n+1$ racines distinctes.
		\item Soit $m\geqslant 3$ entier et $(A,B,C)\in\C[X]^3$ un triplet de polynômes premiers entre eux dans leur ensemble tel que $A^m+B^m=C^m$. Montrer que $A$, $B$ et $C$ sont tous constants.
	\end{enumerate}
\end{pro}

\section{Exercices}

\begin{ex}[Critère d'Eisenstein]
	Soit $p$ un nombre premier et $P\in\Z[X]$ qu'on écrit sous la forme $$P=\sum_{k=0}^{n}a_kX^k$$ Supposons que l'on ait :
	\begin{itemize}
		\item $\forall k\in\llbracket0,n-1\rrbracket,\ p\mid a_k$ ;
		\item $p^2\nmid a_0$.
		\item $p\nmid a_n$
	\end{itemize}
	Montrer que $P$ est irréductible dans $\Z[X]$.
\end{ex}

\begin{ex}[Polynômes à valeurs dans les nombres premiers]
	On note $\mathcal{P}$ l'ensemble des nombres premiers. Trouver tous les polynômes $A \in \Z[X]$ tels que $\forall n \in \N, \ A(n) \in \mathcal{P}$.
\end{ex}

\chapter{Révisions sur les espaces vectoriels}

\section{Problèmes}

\section{Exercices}

\begin{ex}
	Soit $n\in\N\st$. On se donne $X_0,\ldots,X_n$ des parties non vides de $\llbracket1,n\rrbracket$. Montrer qu'il existe $I$ et $J$ deux parties disjointes non vides de $\llbracket0,n\rrbracket$ telles que : $$\bigcup_{i\in I} X_i=\bigcup_{i\in J} X_j$$
\end{ex}

\begin{ex}
	Soit $E$ un $\R$-espace vectoriel et $(e_i)_{i\in I}$ une famille de vecteurs de $E$. On dit que $(e_i)_{i\in I}$ est positivement génératrice lorsque tout vecteur de $E$ est combinaison linéaire à coefficients strictement positifs des $e_i$. Dans tout l'exercice, on suppose $E$ de dimension finie $n\in\N\st$ et on se donne $(e_1,\ldots,e_p)\in E^p$ une famille positivement génératrice.
	\begin{enumerate}
		\item Montrer : $p\geqslant n+1$.
		\item Donner un exemple où on peut avoir $p=n+1$.
		\item On suppose de plus $p\geqslant 2n+1$. Montrer que $(e_1,\ldots,e_p)$ possède une sous-famille stricte qui est toujours positivement génératrice.
		\item Donner un exemple où $p=2n$ et où aucune sous-famille stricte n'est positivement génératrice.
	\end{enumerate}
\end{ex}

\begin{ex}
	Soit $E$ un espace vectoriel de dimension finie, $G$ un sous-groupe fini de $\mathcal{GL}(E)$ et $F$ un sous-espace vectoriel de $E$ stable par tous les éléments de $G$. Montrer que $F$ admet un supplémentaire stable par tous les éléments de $G$.
\end{ex}

\chapter{Révisions sur les matrices}

\section{Problèmes}

\section{Exercices}

\begin{ex}
	Soit $A\in\mathcal{M}_n(\K)$ non scalaire. Montrer que pour tout $(a_1,\ldots,a_n)\in\K^n$ tel que $\Tr(A)=a_1+\cdots+a_n$, il existe une matrice semblable à $A$ dont les coefficients diagonaux sont $a_1,\ldots,a_n$.
\end{ex}

\begin{ex}
	Soit $n\in\N\st$. Déterminer le $k\in\N$ maximal tel qu'il existe des parties $E_1,\ldots,E_k$ de $\llbracket 1,n\rrbracket$ telles que :
	\begin{itemize}
		\item $\card(E_i)$ soit impair pour tout $i$ ;
		\item $\card(E_i\cap E_j)$ soit pair pour tout $i\ne j$.
	\end{itemize}
\end{ex}

\begin{ex}
	Soit $n\geqslant1$. Pour tout $\sigma\in\mathfrak{S}_n$, on note $\text{inv}(\sigma)$ le nombre de points de $\llbracket1,n\rrbracket$ invariants par $\sigma$. Calculer : $$\sum_{\sigma\in\mathfrak{S}_n}\frac{\varepsilon(\sigma)}{\text{inv}(\sigma)+1}$$
\end{ex}

\chapter{Suites et séries numériques}

\section{Problèmes}

\begin{pro}
	Soit $\displaystyle\sum_{n\geqslant 1}a_n$ une série à termes positifs convergente, de somme notée $A$. On pose pour tout $n\in\N\st$ : $u_n=\sqrt[n]{a_1\cdots a_n}$.
	\begin{enumerate}
		\item Démontrer l'inégalité arithmético-géométrique.
		\item Montrer que $\displaystyle\sum_{n\geqslant 1}u_n$ converge et qu'en notant $U$ sa somme, on a $U\leqslant e A$.
		\item Montrer que si une constante $k$ vérifie l'inégalité $U\leqslant k A$ pour toute série à termes positifs convergente $\displaystyle\sum_{n\geqslant 1}a_n$, alors $e\leqslant k$.
	\end{enumerate}
\end{pro}

\begin{pro}
	Convergence et somme de : $$\sum_{n\geqslant 2}\frac{(-1)^n}{n}\left\lfloor\frac{\ln(n)}{\ln(2)}\right\rfloor$$
\end{pro}

\section{Exercices}

\begin{ex}
	Soit $u_0$ un réel strictement positif. On définit par récurrence : $$\forall n\in\N,\ u_{n+1}=u_n^{u_n}$$
	\begin{enumerate}
		\item Etudier la convergence de la suite $(u_n)$.
		\item Désormais et jusqu'à la fin de l'exercice, on suppose $u_0\in\N\setminus\{0,1\}$. Montrer : $$\forall N\in\N,\ \forall n\in\llbracket0,N\rrbracket,\ u_n\mid u_N$$
		\item Montrer que pour tous $N$ et $k$ dans $\N$, on a : $u_{N+k}\geqslant u_N^{k+1}$.
		\item Montrer la convergence de la série de terme général $1/u_n$. On note $S$ sa somme.
		\item Montrer : $$u_N\sum_{n=N+1}^{+\infty}\frac{1}{u_n}\xrightarrow[N\to+\infty]{}0$$
		\item Montrer que $S\notin\Q$.
	\end{enumerate}
\end{ex}

\begin{ex}
	Pour $x\in\R$, étudier la convergence de la suite de terme général $P_n(x)=\displaystyle\prod_{k=1}^{n}\cos\left(\frac{x}{2^k}\right)$.
\end{ex}

\begin{ex}
	On définit la suite $(u_n)_{n\in\N}$ par $u_0=1$ et $$\forall n\in\N,\ u_{n+1}=e^{-u_n}-1$$ Déterminer un équivalent de la suite $(u_n)_{n\in\N}$.
\end{ex}

\chapter{Réduction des endomorphismes}

\section{Problèmes}

\section{Exercices}

\begin{ex}
	Soit $A\in\mathcal{M}_n(\Z/p\Z)$ où $n\in\N\st$ et $p$ est un nombre premier. Montrer que $A$ est diagonalisable si, et seulement si, $A^p=A$.
\end{ex}

\begin{ex}
	Pour $\sigma\in\mathcal{S}_n$, on note $P_\sigma\in\mathcal{M}_n(\C)$ la matrice de permutation associée à $\sigma$, de terme général $\delta_{i,\sigma(j)}$. Soit $(\sigma,\tau)\in\mathcal{S}_n^2$. Montrer que $\sigma$ et $\tau$ sont conjuguées dans le groupe $\mathcal{S}_n$ si, et seulement si, les matrices $P_\sigma$ et $P_\tau$ sont semblables.
\end{ex}

\begin{ex}
	En considérant des sous-groupes dont tous les éléments sont des matrices de symétries, montrer que si $\mathcal{GL}_n(\C)$ et $\mathcal{GL}_m(\C)$ sont isomorphes, alors $n=m$.
\end{ex}

\begin{ex}
	On admet le théorème de réduction de Jordan. Montrer que toute matrice inversible à coefficients complexes admet une racine carrée. Démontrer le théorème de réduction de Jordan.
\end{ex}

\chapter{Espaces euclidiens}

\section{Problèmes}

\begin{pro}[Autour du théorème de Schur-Horn]
	
\end{pro}

\begin{pro}[Autour des égalités de Courant-Fischer]
	On considère $u\in S(E)$ avec $E$ euclidien de dimension $n$. On note $\lambda_1\leqslant\ldots\leqslant\lambda_n$ ses valeurs propres. On note $\mathcal{F}_k$ l'ensemble des sous-espaces vectoriels de $E$ de dimension $k$ et $S$ la sphère unité de $E$.
	\begin{enumerate}
		\item Montrer que : $$\lambda_k=\inf_{F\in\mathcal{F}_k}\sup_{x\in S\cap F}(x|u(x))$$ puis montrer que $$\lambda_{n+1-k}=\sup_{F\in\mathcal{F}_k}\inf_{x\in S\cap F}(x|u(x))$$
		\item Soit $A\in\mathcal{S}_n(\R)$ avec $n\geqslant 2$. On note $B$ la matrice extraite de $A$ en lui supprimant sa dernière ligne et sa dernière colonne. On note $\lambda_1\leqslant\ldots\leqslant\lambda_n$ les valeurs propres de $A$ et $\mu_1\leqslant\ldots\leqslant\mu_{n-1}$ celles de $B$. Montrer : $$\lambda_1\leqslant\mu_1\leqslant\lambda_2\leqslant\ldots\leqslant\mu_{n-1}\leqslant\lambda_n$$
		\item Pour $A\in\mathcal{S}_n(\R)$, on note $\lambda_1(A)\geqslant\ldots\geqslant\lambda_n(A)$ les valeurs propres de $A$, comptées avec multiplicité. Montrer : $$\forall (A,B)\in\mathcal{S}_n(\R)^2,\ \lambda_{i+j-1}(A+B)\leqslant\lambda_i(A)+\lambda_j(B)$$
	\end{enumerate}
\end{pro}

\begin{pro}[Autour des racines carrées]
	On présente ici la racine carré d'un endomorphisme symétrique positif et quelques applications.
	\begin{enumerate}
		\item Soit $E$ un espace euclidien ainsi que $u\in\mathcal{S}^+(E)$. Montrer qu'il existe un unique endomorphisme $v\in\mathcal{S}^+(E)$ tel que $u=v^2$. Montrer que $v$ est un polynôme en $u$.
		\item Soit $n\in\N\st$ et $M\in\mathcal{GL}_n(\R)$. Montrer qu'il existe un unique couple $(O,S)\in\mathcal{O}_n(\R)\times\mathcal{S}_n^{++}(\R)$ tel que $M=OS$. Montrer que l'on conserve l'existence si l'on suppose seulement $M\in\mathcal{M}_n(\R)$. Montrer que l'application qui à $(O,S)\in\mathcal{O}_n(\R)\times\mathcal{S}_n^{++}(\R)$ associe $OS$ est un homéomorphisme de $\mathcal{O}_n(\R)\times\mathcal{S}_n^{++}(\R)$ sur $\mathcal{GL}_n(\R)$.
		\item Soit $A\in\mathcal{S}_n^{++}(\R)$ et $B\in\mathcal{S}_n(\R)$. Montrer qu'il existe $P\in\mathcal{GL}_n(\R)$ telle que $P^TAP=I_n$ et $P^TBP$ soit diagonale.
	\end{enumerate}
\end{pro}

\section{Exercices}

\begin{ex}
	On projette orthogonalement un cube de l'espace de côté $a$ sur un plan. On note $x$, $y$ et $z$ les longueurs des arêtes du projeté. Montrer que $x^2+y^2+z^2=2a^2$.
\end{ex}

\begin{ex}
	Soit $E$ un espace euclidien de dimension $n$. On suppose qu'il existe $(x_1,\ldots,x_p)\in E^{p}$ telle que : $$\forall (i,j)\in\llbracket1,p\rrbracket^2,\ i\ne j\implies (x_i\mid x_j)<0$$ Montrer que $p\leqslant n+1$.
\end{ex}

\begin{ex}
	Soit $E$ un espace euclidien. Montrer que le groupe $\mathcal{O}(E)$ des isométries vectorielles est engendré par l'ensemble des réflexions.
\end{ex}

\begin{ex}
	Soit $n\in\N\st$ et $M\in\mathcal{M}_n(\R)$. On suppose qu'il existe un entier $k\geqslant 2$ tel que $M^k=M^T$.
	\begin{enumerate}
		\item Montrer que $M^TM$ est la matrice d'une projection orthogonale.
		\item Montrer que $M$ est $\C$-diagonalisable.
	\end{enumerate}
\end{ex}

\begin{ex}
	Soit $(A,B)\in\mathcal{S}_n(\R)^2$ vérifiant : $$\forall X\in\R^n,\ 0\leqslant X^TAX\leqslant X^TBX$$ Montrer : $\det(A)\leqslant\det(B)$.
\end{ex}

\begin{ex}
	Soit $S\in\mathcal{S}_n^{++}(\R)$ et $A\in\mathcal{A}_n(\R)$. Montrer que $SA$ est diagonalisable sur $\C$ et à spectre inclus dans $i\R$.
\end{ex}

\chapter{Topologie}

\section{Problèmes}

\begin{pro}[Autour de la propriété de Borel-Lebesgue]
	On se donne $K$ un compact.
	\begin{enumerate}
		\item  Précompacité : montrer que pour tout $\varepsilon>0$, il existe une partie finie $C$ de $K$ telle que $$K\subset\bigcup_{x\in C}B(x,\varepsilon)$$
		\item Propriété de Borel-Lebesgue :
		\begin{enumerate}[label=\alph*.]
			\item Soit $(O_i)_{i\in I}$ une famille d'ouverts de $K$ recouvrant $K$. Montrer qu'il existe $\varepsilon>0$ tel que $$\forall x\in K,\ \exists i\in I\ : \ B_o(x,\varepsilon)\subset O_i$$
			\item Montrer que pour toute famille $(O_i)_{i\in I}$ d'ouverts de $A$ recouvrant $K$, il existe une partie finie $J$ de $I$ telle que $(O_j)_{j\in J}$ recouvre $K$. On pourra utiliser la précompacité de $K$.
			\item Montrer que la propriété de la deuxième question est suffisante pour que $K$ soit compacte. On pourra utiliser le fait que l'ensemble des valeurs d'adhérence d'une suite $(u_n)$ est $\displaystyle\bigcap_{n\in\N}\overline{\{u_k\mid k\geqslant n\}}$
		\end{enumerate}
		\item Applications :
		\begin{enumerate}[label=\alph*.]
			\item Montrer que les idéaux maximaux de $\mathcal{C}^0([0,1],\R)$ sont de la forme : $$I_x=\left\{f\in\mathcal{C}^0([0,1],\R)\ \mid\ f(x)=0\right\}$$
			\item Soit $E$ un espace métrique compact et $f:E\to\R$ une fonction localement bornée, c'est-à-dire que pour tout $x\in E$, il existe un voisinage $V_x$ de $x$ tel que pour tout $y\in V_y$, on ait $\lvert f(y)\rvert\leqslant M_x$. Montrer que $f$ est bornée.
			\item On admet le théorème de Kakutani : pour tout convexe compact non vide $K$ d'un espace vectoriel normé $E$ et toute application affine continue $u$ qui stabilise $K$, $u$ possède un point fixe. Montrer qu'une famille $(u_i)_{i\in I}$ d'applications affines qui commutent deux à deux et stabilisent toutes un compact convexe non vide $K$ possèdent un point fixe commun.
		\end{enumerate}
	\end{enumerate}
\end{pro}

\begin{pro}[Autour du théorème de Baire]
	Soit $A$ une partie localement compacte d'un espace vectoriel normé (\textit{ie} dans laquelle tout point possède un voisinage compact).
	\begin{enumerate}
		\item Soit $(O_n)$ une suite d'ouverts de $A$, tous denses dans $A$. Montrer que $$\Omega=\bigcap_{n\in\N} O_n$$ est dense dans $A$. Est-il nécessairement ouvert ?
		\item En déduire que si une suite $(F_n)$ de fermés de $A$ a une réunion d'intérieur non vide, alors l'un des $F_n$ est d'intérieur non vide.
	\end{enumerate}
\end{pro}

\begin{pro}[Connexités par arcs]
	\begin{enumerate}
		\item Montrer que l'ensemble $\mathcal{GL}_n(\C)$ est connexe par arcs.
		\item Montrer que $\mathcal{GL}_n(\R)$ n'est pas connexe par arcs et déterminer ses composantes connexes par arcs.
		\item Déterminer les composantes connexes par arcs de l'ensemble des matrices de projecteurs pour $\K=\C$ puis pour $\K=\R$.
	\end{enumerate}
\end{pro}

\section{Exercices}

\begin{ex}[Point fixe de Kakutani]
	Soit $K$ un compact convexe non vide et $u$ une application affine continue qui stabilise $K$. Montrer que $u$ admet un point fixe.
\end{ex}

\begin{ex}
	Déterminer les sous-groupes fermés de $(\C\st,\times)$.
\end{ex}

\begin{ex}
	Montrer que pour tout polynôme $P\in\R[X]$, il existe une norme sur $\R[X]$ qui fait converger la suite $(X^k)_{k\in\N}$ vers $P$.
\end{ex}

\begin{ex}
	Montrer qu'une matrice $A\in\mathcal{M}_n(\K)$ est nilpotente si, et seulement si, la matrice nulle est adhérente à la classe de similitude de $A$.
\end{ex}

\begin{ex}
	Montrer qu'il n'existe pas de norme sur $\mathcal{C}^{\infty}([0,1],\R)$ pour laquelle la dérivation soit continue.
\end{ex}

\begin{ex}
	Soit $C$ une partie convexe compacte non vide et $f:C\to C$ 1-lipschitzienne. Montrer que $f$ admet au moins un point fixe. On pourra commencer par le cas où $f$ est $k$-lipschitzienne avec $k<1$.
\end{ex}

\begin{ex}
	Soit $A$ une partie non triviale d'un espace vectoriel normé $E$. Montrer que la frontière de $A$ est non vide.
\end{ex}

\begin{ex}
	Soit $G$ un sous-groupe additif de $\R^n$. On suppose que $G$ est fermé et que pour tout $y\in G$ et pour tout $\varepsilon>0$, il existe $x\in G$ tel que $0<\lVert x-y\rVert<\varepsilon$. Montrer que $G$ contient une demi-droite.
\end{ex}

\begin{ex}
	Soit $(E,d)$ un espace métrique et $K$ une partie compacte de $E$. On considère $f$ une fonction de $K$ dans $K$ telle que : $$\forall (x,y)\in K^2,\ d(f(x),f(y))\geqslant d(x,y)$$
	Montrer que $f(K)$ est dense dans $K$ et en déduire que $f$ est un homéomorphisme.
\end{ex}

\begin{ex}[Lemme d'Auerbach]
	Soit $(E,\lVert\cdot\rVert)$ un espace vectoriel normé de dimension finie. Montrer qu'il existe une base de $E$ unitaire dont la base duale est unitaire (pour la norme subordonnée).
\end{ex}

\chapter{Intégrales généralisées}

\begin{ex}
	Pour $n\in\N$ et $x\in\R$, on pose : $H_n(x)=\displaystyle\frac{1}{n!}\int_{-x}^{x}(x^2-t^2)^n\cos(t)\ dt$
	\begin{enumerate}
		\item Montrer que pour tout $n\in\N$, il existe des polynômes $C_n$ et $S_n$ à coefficients entiers de degrés inférieurs ou égaux à $n$ tels que pour tout $x\in\R$, on ait : $H_n(x)=C_n(x)\cos(x)+S_n(x)\sin(x)$.
		\item Soit $\theta\in\Q\st$. Montrer que $\tan(\theta)$ est irrationnel.
	\end{enumerate}
\end{ex}

\begin{ex}
	Calculer $\displaystyle\int_{0}^{\pi/2}\ln(\sin(x))dx$.
\end{ex}

\chapter{Intégrales à paramètres}

\section{Problèmes}

\begin{pro}
	Pour $(x,y)\in\R_+\st\times\R$, on pose $$G(x,y)=\int_{0}^{+\infty}\frac{e^{-xt}e^{iyt}}{\sqrt{t}}\ dt$$
	\begin{enumerate}
		\item Montrer que $G$ est continue sur $\R_+\st\times\R$.
		\item Calculer $G(1,0)$.
		\item Exprimer $G(x,y)$ à l'aide de $G(1,z)$ pour un $z$ bien choisi.
		\item Montrer que $H:z\mapsto G(1,z)$ est dérivable sur $\R$ et que $$\forall z\in\R,\ H'(z)=\frac{i}{2(1-iz)}H(z)$$
		\item En déduire la valeur de $H(z)$ pour tout $z\in\R$ puis celle de $G(x,y)$ pour tout $(x,y)\in\R_+\st\times\R$.
		\item Montrer que $G(0,y)$ a un sens pour tout $y>0$.
		\item Montrer que $\displaystyle\lim_{x\to 0^+}G(x,1)=G(0,1)$. En déduire la convergence et la valeur de $$I=\int_{0}^{+\infty}\cos(t^2)\ dt \ \ \text{et} \ \ J=\int_{0}^{+\infty}\sin(t^2)\ dt$$
	\end{enumerate}
\end{pro}

\begin{pro}
	Soit $f$ une fonction continue et à valeurs strictement positives sur $[0,a]$ avec $a>0$. on suppose qu'il existe $(\alpha,p)\in(\R_+\st)^2$ tel que $f(t)=1-\alpha t^p+o(t^p)$ au voisinage de $0$ et $\forall t\in]0,a],\ f(t)<1$. On se propose de trouver un équivalent de $$I(x)=\int_{0}^{a}f(t)^x\ dt$$ lorsque $x$ tend vers $+\infty$.
	\begin{enumerate}
		\item Montrer qu'il existe $\beta>0$ tel que $$\forall t\in[0,a],\ f(t)\leqslant\exp(-\beta t^p)$$
		\item Déterminer la limite, quand $x$ tend vers $+\infty$, de $$g(x)=\int_{0}^{ax^{1/p}}f\left(ux^{-1/p}\right)^x\ dt$$
		\item Conclure.
		\item Application : convergence, limite et équivalent lorsque $n$ tend vers l'infini de la suite $(I_n)$ définie par $$I_n=\int_{0}^{+\infty}\frac{dt}{(1+t^3)^n}$$
	\end{enumerate}
\end{pro}

\begin{pro}
	Si $f$ et $g$ sont des fonctions continues par morceaux de $\R$ dans $\R$, on appelle produit de convolution de $f$ et $g$, $f\star g$, la fonction, si elle existe, définie par $$f\star g:x\mapsto\displaystyle\int_\R f(x-t)g(t)\ dt$$
	\begin{enumerate}
		\item Montrer que si $f$ et $g$ sont à support compact, alors $f\star g$ est aussi à support compact.
		\item Démontrer que le produit de convolution est commutatif. On admet qu'il est associatif.
		\item Montrer que si $f$ est bornée et $g$ est intégrable, alors $f\star g$ existe et est bornée par $\lVert f\rVert_\infty\lVert g\rVert_1$.
		\item Pour $a>0$, on note $\varphi_a$ la fonction $t\mapsto\exp(-a^2t^2)$. Montrer que $$\varphi_a\star\varphi_b=\sqrt{\frac{\pi}{a^2+b^2}}\varphi_{\frac{ab}{\sqrt{a^2+b^2}}}$$
		\item Soit $f$ continue bornée sur $\R$. Soit $(\varphi_n)$ une suite de fonctions continues positives intégrables telle que $\displaystyle\int_\R\varphi_n(t)\ dt\xrightarrow[n\to+\infty]{}1$ et, pour tout $\delta>0$, $\displaystyle\int_{\lvert t\rvert\geqslant \delta}\varphi_n(t)\ dt\xrightarrow[n\to+\infty]{}0$. Montrer que la suite de fonctions $(f\star\varphi_n)$ converge simplement vers $f$, et uniformément sur tout segment.
		\item Pour $n\in\N$, on pose $\displaystyle g_n(x)=\frac{1}{c_n}(1-x^2)^n$ si $x\in[-1,1]$ avec $c_n=\displaystyle\int_{-1}^{1}(1-t^2)^n\ dt$. Soit $f:[-1,1]\to\R$ nulle en dehors de $\displaystyle\left[-\frac{1}{2},\frac{1}{2}\right]$. Montrer que la restriction de $f\star g_n$ à $\displaystyle\left[-\frac{1}{4},\frac{1}{4}\right]$ est polynomiale.Montrer que la suite de fonctions $(f\star g_n)$ converge uniformément vers $f$ sur $\displaystyle\left[-\frac{1}{4},\frac{1}{4}\right]$ et en déduire une preuve du théorème de Weierstrass.
	\end{enumerate}
\end{pro}

\section{Exercices}

\begin{ex}[Démonstration du théorème de D'Alembert-Gauss]
	Soit $f$ une fonction de $\R$ dans $\C\st$, $2\pi$-périodique et de classe $\mathcal{C}^1$.
	\begin{enumerate}
		\item Montrer : $$\frac{1}{2i\pi}\int_{0}^{2\pi}\frac{f'(t)}{f(t)}dt\in\Z$$
		\item En déduire le théorème de D'Alembert-Gauss.
		\item Calculer la valeur de cette intégrale pour $P\in\C[X]$.
	\end{enumerate}
\end{ex}

\begin{ex}
	Soit $a$ et $b$ deux réels strictement supérieurs à $1$. Calculer : $$\int_{0}^{\pi}\ln\left(\frac{b-\cos(x)}{a-\cos(x)}\right)dx$$
\end{ex}

\begin{ex}[Formule sommatoire de Poisson]
	On note $\mathcal{S}$ l'ensemble des fonctions de $\R$ dans $\C$ de classe $\mathcal{C}^\infty$ telles que pour tous entiers naturels $n$ et $m$, on ait $x^mf^{(n)}(x)\xrightarrow[x\to\pm\infty]{}0$. Soit $f\in\mathcal{S}$. On pose : $$\forall x\in \R,\ \varphi(x)=\sum_{k\in\Z}f(x+k)$$
	\begin{enumerate}
		\item Bonne définition de $\varphi$ ? Régularité de $\varphi$ ? Une autre remarque ?
		\item Pour $n\in\Z$, calculer : $$\int_{0}^{1}\varphi(t)e^{-2i\pi nt}dt$$
		\item On admet qu'une fonction $\psi$ de classe $\mathcal{C}^\infty$ et $1$-périodique s'écrit sous la forme : $$\psi(x)=\sum_{n\in\Z}c_ne^{2i\pi nx}$$ On pose : $$\forall n\in\Z,\ \hat{f}(n)=\int_{\R}f(t)e^{-2i\pi nt}dt$$ Montrer la formule sommatoire de Poisson : $$\sum_{n\in\Z}f(n)=\sum_{n\in\Z}\hat{f}(n)$$
	\end{enumerate}
\end{ex}

\chapter{Équations différentielles}

\section{Problèmes}

\section{Exercices}

\begin{ex}
	Soit $f\in\mathcal{C}^\infty(\R,\R)$ telle que l'espace vectoriel $\text{Vect}(x\mapsto f(x+a))_{a\in\R}$ soit de dimension finie. Montrer que $f$ est solution d'une équation différentielle linéaire à coefficients constants.
\end{ex}

\chapter{Algèbre générale}

\section{Problèmes}

\begin{pro}[Théorème de Wedderburn]
	
\end{pro}

\section{Exercices}


\chapter{}

\section{Problèmes}

\section{Exercices}

\end{document}